%%%% ABSTRACT
%%
%% Versão do resumo para idioma de divulgação internacional.

\begin{abstractutfpr}%% Ambiente abstractutfpr

Image fusion for panoramic mosaic construction is a fundamental technique in computer vision with applications in surveillance, autonomous navigation, and environmental monitoring. Despite recent advances with deep learning techniques, the literature lacks a systematic comparison between classical and modern approaches considering multiple quality and efficiency criteria. Additionally, the impact of fusion quality on subsequent tasks such as object detection remains underexplored. This work presents a comprehensive comparative analysis of 20 image fusion techniques, 15 based on classical methods and 5 on deep learning, applied to the Robust Image Stitching with Multiple Registrations dataset. 295 successful executions were performed on 15 image sets, with quantitative evaluation through PSNR, SSIM, MSE, Ghosting Score, Seam Energy, DeltaE, Entropy, Sharpness, and processing time metrics. Results show that OpenCV Stitcher and Basic Stitcher methods achieved the best overall performance (score 0.8732), with PSNR of 11.10 dB and average time of 0.26 seconds, representing excellent balance between quality and efficiency. The Seams Exposure Multiband method achieved the highest PSNR (11.49 dB) and lowest seam energy (77.42), indicating better visual quality, but with significantly higher computational cost (7.11s). The comparison between classical and deep learning approaches revealed a small quality difference (+0.08 dB PSNR, +0.0066 SSIM for deep learning), but with 68\% higher average processing time for neural network-based techniques. All methods presented high Ghosting Score (~0.94), reflecting the inherent challenges of the dataset used. In the second part of the work, the impact of different fusion techniques on object detection is investigated by applying YOLO models (YOLOv5, YOLOv8, and YOLO-NAS) to the generated panoramas and comparing with detections in individual images. \textcolor{blue}{[Detection results will be added after completion of experiments.]} The results demonstrate that classical techniques offer excellent cost-benefit ratio for panoramic fusion, being suitable for applications requiring fast processing, while deep learning methods are indicated when visual quality is an absolute priority. The work provides practical guidelines for choosing fusion techniques considering the application context and establishes a methodological basis for future investigations on the integration of image fusion and object detection.

\vspace{1em}

\noindent\textbf{Keywords}: image fusion; panoramic mosaic; image stitching; computer vision; deep learning; object detection; YOLO.

\end{abstractutfpr}
